\begin{frame}[noprogressbar,noframenumbering]{Vì sao fact\symbol{95}order\symbol{95}fulfillment là accumulating snapshot? (3)}
\begin{itemize}
  \item Đơn đi qua nhiều mốc; một dòng, cập nhật theo mốc; lead time tính sẵn (approval 10,3 giờ, handling 2,9 ngày, transit 9,3 ngày)
  \item Slide chính liên quan: S63
\end{itemize}
\begin{center}\adjustbox{max width=\linewidth,max height=0.5\textheight}{\scriptsize\renewcommand{\arraystretch}{1.15}\begin{tabular}{ll}\toprule
\textbf{Cột của fact\_order\_fulfillment (đơn 19c36691…)} & \textbf{Giá trị} \\\midrule
approval\_hours & 36,61 \\
handling\_days & 0,71 \\
transit\_days & 12,84 \\
lead\_days & 15,07 \\
promised\_days & 21 \\
delay\_days & -6 (giao sớm 6 ngày) \\
order\_value & 43,00 \\
freight\_total & 45,30 \\
payment\_total & 88,30 \\
\bottomrule\end{tabular}}\end{center}

\note{[S63] B01 - dự phòng cho câu hỏi: Vì sao \texttt{fact\_\allowbreak{}order\_\allowbreak{}fulfillment} là accumulating snapshot? (3). Xem huong-dan/04-cau-hoi-van-dap.md.}
\end{frame}

\begin{frame}[noprogressbar,noframenumbering]{Vì sao chọn ROLAP mà không MOLAP? (25)}
\begin{itemize}
  \item PostgreSQL không có khối đa chiều, dữ liệu nhỏ, kho luôn tươi, Power BI là lớp ngữ nghĩa
  \item Slide chính liên quan: S11, S87
\end{itemize}
\begin{center}\adjustbox{max width=\linewidth,max height=0.5\textheight}{\scriptsize\renewcommand{\arraystretch}{1.15}\begin{tabular}{llll}\toprule
\textbf{Tiêu chí} & \textbf{MOLAP} & \textbf{ROLAP} & \textbf{HOLAP} \\\midrule
Lưu trữ & khối đa chiều riêng & bảng quan hệ (lược đồ sao) & kết hợp hai cách \\
Tốc độ truy vấn & rất nhanh (tính sẵn) & phụ thuộc chỉ mục, materialized view & cân bằng \\
Độ tươi dữ liệu & phải xử lý lại khối & luôn tươi sau mỗi lần nạp & tùy phần \\
Công cụ & ví dụ SSAS & PostgreSQL + Power BI & -- \\
Nhóm chọn & -- & \ding{51} ROLAP & -- \\
\bottomrule\end{tabular}}\end{center}

\note{[S11] B02 - dự phòng cho câu hỏi: Vì sao chọn ROLAP mà không MOLAP? (25). Xem huong-dan/04-cau-hoi-van-dap.md.}
\end{frame}

\begin{frame}[noprogressbar,noframenumbering]{SCD2 hoạt động thế nào? (6)}
\begin{itemize}
  \item Mỗi lần zip đổi mở phiên bản mới; 96.096 khách $\rightarrow$ 96.339 dòng; fact tra theo ngày mua
  \item Slide chính liên quan: S59, S75
\end{itemize}
\begin{center}\adjustbox{max width=\linewidth,max height=0.5\textheight}{\scriptsize\renewcommand{\arraystretch}{1.15}\begin{tabular}{llllllll}\toprule
\textbf{customer\_key} & \textbf{khách thật} & \textbf{zip} & \textbf{thành phố} & \textbf{effective\_from} & \textbf{effective\_to} & \textbf{is\_current} & \textbf{version} \\\midrule
15081 & 570eb70f… & 31930 & belo horizonte & 1900-01-01 & 2017-11-08 & không & 1 \\
15082 & 570eb70f… & 31130 & belo horizonte & 2017-11-09 & 9999-12-31 & có & 2 \\
\bottomrule\end{tabular}}\end{center}

\note{[S59] B03 - dự phòng cho câu hỏi: SCD2 hoạt động thế nào? (6). Xem huong-dan/04-cau-hoi-van-dap.md.}
\end{frame}

\begin{frame}[noprogressbar,noframenumbering]{Xử lý 1 triệu dòng geolocation ra sao? (20)}
\begin{columns}[T]\begin{column}{0.4\textwidth}
\begin{itemize}
  \item Lọc khung Brazil (loại 42), trung vị toạ độ, mode thành phố, 19.015 zip; zip chỉ có toạ độ ngoài khung vẫn có dòng
  \item Slide chính liên quan: S40, S72
\end{itemize}
\end{column}\begin{column}{0.58\textwidth}
\begin{center}\anh[0.74\textheight]{3-4-3_before_after_geolocation.png}\end{center}
\end{column}\end{columns}

\note{[B04] B04 - dự phòng cho câu hỏi: Xử lý 1 triệu dòng geolocation ra sao? (20). Xem huong-dan/04-cau-hoi-van-dap.md.}
\end{frame}

\begin{frame}[noprogressbar,noframenumbering]{Vì sao đúng hạn 93,2\% mà review đơn trễ vẫn tệ? (33)}
\begin{columns}[T]\begin{column}{0.4\textwidth}
\begin{itemize}
  \item Trễ chỉ 6,8\% đơn nhưng tạo 27,4\% review tiêu cực; điểm 4,29 $\rightarrow$ 1,70; ngày dự kiến quá rộng
  \item Slide chính liên quan: S45, S98
\end{itemize}
\end{column}\begin{column}{0.58\textwidth}
\begin{center}\anh[0.74\textheight]{3-3-1_delay_vs_review.png}\end{center}
\end{column}\end{columns}

\note{[B05] B05 - dự phòng cho câu hỏi: Vì sao đúng hạn 93,2\% mà review đơn trễ vẫn tệ? (33). Xem huong-dan/04-cau-hoi-van-dap.md.}
\end{frame}

\begin{frame}[noprogressbar,noframenumbering]{Vì sao star schema chỉ nhanh hơn OLTP 1,0–1,6 lần mà vẫn đáng dùng? (30)}
\begin{columns}[T]\begin{column}{0.4\textwidth}
\begin{itemize}
  \item Dữ liệu nhỏ; giá trị là truy vấn đơn giản, KPI nhất quán; tốc độ từ materialized view và index
  \item Slide chính liên quan: S88
\end{itemize}
\end{column}\begin{column}{0.58\textwidth}
\begin{center}\anh[0.74\textheight]{3-5-1_benchmark.png}\end{center}
\end{column}\end{columns}

\note{[B06] B06 - dự phòng cho câu hỏi: Vì sao star schema chỉ nhanh hơn OLTP 1,0–1,6 lần mà vẫn đáng dùng? (30). Xem huong-dan/04-cau-hoi-van-dap.md.}
\end{frame}

\begin{frame}[noprogressbar,noframenumbering]{Làm sao biết nạp tăng dần đúng? (13)}
\begin{itemize}
  \item So với nạp full: 16/16 bảng khớp số dòng, tổng tiền, checksum; chạy lại một đợt cho 0 dòng mới
  \item Slide chính liên quan: S80
\end{itemize}


\note{[B07] B07 - dự phòng cho câu hỏi: Làm sao biết nạp tăng dần đúng? (13). Xem huong-dan/04-cau-hoi-van-dap.md.}
\end{frame}

\begin{frame}[noprogressbar,noframenumbering]{GMV khác tổng giá hàng thế nào? (23)}
\begin{itemize}
  \item GMV loại đơn Failed: 13.494.400,74 so với 13.591.643,70
  \item Slide chính liên quan: S26, S30
\end{itemize}


\note{[B08] B08 - dự phòng cho câu hỏi: GMV khác tổng giá hàng thế nào? (23). Xem huong-dan/04-cau-hoi-van-dap.md.}
\end{frame}

\begin{frame}[noprogressbar,noframenumbering]{Dòng mồ côi được xử lý ra sao? (15)}
\begin{itemize}
  \item Loại kèm lý do, thử lại ở đợt sau; mới chứng minh bằng test
  \item Slide chính liên quan: S81
\end{itemize}


\note{[B09] B09 - dự phòng cho câu hỏi: Dòng mồ côi được xử lý ra sao? (15). Xem huong-dan/04-cau-hoi-van-dap.md.}
\end{frame}

\begin{frame}[noprogressbar,noframenumbering]{Ngày hứa giao so với thực tế cho thấy điều gì? (34)}
\begin{itemize}
  \item Hứa 24,4, thực tế 12,6: sớm 11,8 ngày; 91,9\% đơn giao sớm
  \item Slide chính liên quan: S27, S98
\end{itemize}


\note{[B10] B10 - dự phòng cho câu hỏi: Ngày hứa giao so với thực tế cho thấy điều gì? (34). Xem huong-dan/04-cau-hoi-van-dap.md.}
\end{frame}

% --- slide lý thuyết rút khỏi bài trình bày chính (barem không có mục lý thuyết)
\begin{frame}[noprogressbar,noframenumbering]{Kimball xây từ các data mart có chiều dùng chung, Inmon xây kho 3NF rồi mới tách mart}
\begin{itemize}
  \item Nhóm chọn hướng Kimball: bus matrix, conformed dimension
\end{itemize}
\begin{center}\adjustbox{max width=\linewidth,max height=0.5\textheight}{\scriptsize\renewcommand{\arraystretch}{1.15}\begin{tabular}{lll}\toprule
\textbf{Tiêu chí} & \textbf{Kimball} & \textbf{Inmon} \\\midrule
Cách xây & từ data mart theo quy trình nghiệp vụ, chiều dùng chung & kho doanh nghiệp chuẩn hóa 3NF, rồi tách data mart \\
Mô hình & chiều (lược đồ sao) & chuẩn hóa \\
Thời gian ra kết quả & nhanh, từng quy trình một & lâu hơn, thiết kế tổng thể trước \\
Nhóm chọn & \ding{51} (bus matrix, \texttt{dim\_seller} dùng chung) & -- \\
\bottomrule\end{tabular}}\end{center}

\note{[S06] Vì bài toán có hai quy trình nghiệp vụ dùng chung chiều seller và ngày, cách Kimball phù hợp và nhanh ra kết quả.}
\end{frame}

\begin{frame}[noprogressbar,noframenumbering]{Mô hình chiều gồm fact, dimension và grain: chọn grain trước, mọi thứ còn lại theo sau}
\begin{columns}[T]\begin{column}{0.4\textwidth}
\begin{itemize}
  \item Grain của dự án: 1 dòng hàng, 1 đơn, 1 lần thanh toán, 1 review, 1 MQL, 1 seller $\times$ tháng
\end{itemize}
\end{column}\begin{column}{0.58\textwidth}
\begin{center}\anh[0.74\textheight]{3-3-4_erd_dw.png}\end{center}
\end{column}\end{columns}

\note{[S07] Grain mịn nhất cho phép cộng lên mọi cấp; ví dụ GMV tính từ grain dòng hàng.}
\end{frame}

\begin{frame}[noprogressbar,noframenumbering]{Ba loại fact và hai kiểu SCD đủ để mô tả mọi nhu cầu của bài toán}
\begin{itemize}
  \item Transaction $\cdot$ periodic snapshot $\cdot$ accumulating snapshot; SCD1 (ghi đè) và SCD2 (giữ lịch sử)
\end{itemize}
\begin{center}\adjustbox{max width=\linewidth,max height=0.5\textheight}{\scriptsize\renewcommand{\arraystretch}{1.15}\begin{tabular}{lll}\toprule
\textbf{Loại} & \textbf{Một dòng là} & \textbf{Ví dụ trong dự án} \\\midrule
Transaction & một sự kiện đã xảy ra & \texttt{fact\_order\_item}, \texttt{fact\_payment}, \texttt{fact\_review} \\
Periodic snapshot & trạng thái của một kỳ & \texttt{fact\_seller\_monthly} (36.583 dòng, dày đặc) \\
Accumulating snapshot & cả vòng đời, cập nhật theo mốc & \texttt{fact\_order\_fulfillment} (99.441), \texttt{fact\_lead\_funnel} (8.000) \\
SCD1 / SCD2 & ghi đè / giữ lịch sử & \texttt{dim\_seller} / \texttt{dim\_customer} \\
\bottomrule\end{tabular}}\end{center}

\note{[S08] Dự án có đủ cả ba loại fact và dùng SCD2 cho khách, SCD1 cho seller.}
\end{frame}

\begin{frame}[noprogressbar,noframenumbering]{ETL biến dữ liệu nguồn bẩn thành dữ liệu đáng tin qua staging và ODS}
\begin{columns}[T]\begin{column}{0.4\textwidth}
\begin{itemize}
  \item Extract $\rightarrow$ làm sạch ở ODS $\rightarrow$ nạp chiều và fact; full và incremental
\end{itemize}
\end{column}\begin{column}{0.58\textwidth}
\begin{center}\anh[0.74\textheight]{3-4-1_pipeline.png}\end{center}
\end{column}\end{columns}

\note{[S09] Nguyên tắc của nhóm là không xóa âm thầm: dòng lỗi được gắn cờ hoặc ghi vào bảng loại kèm lý do.}
\end{frame}

\begin{frame}[noprogressbar,noframenumbering]{Năm thao tác OLAP đều quy về SQL trên lược đồ sao}
\begin{itemize}
  \item Roll-up, drill-down, slice, dice, pivot
\end{itemize}
\begin{center}\adjustbox{max width=\linewidth,max height=0.5\textheight}{\scriptsize\renewcommand{\arraystretch}{1.15}\begin{tabular}{lll}\toprule
\textbf{Thao tác} & \textbf{SQL} & \textbf{Ví dụ trong dự án} \\\midrule
Roll-up & \texttt{ROLLUP} & quý $\rightarrow$ năm $\rightarrow$ tổng \\
Drill-down & thêm cấp vào \texttt{GROUP BY} & 2017 $\rightarrow$ Q4 $\rightarrow$ 11/2017 $\rightarrow$ ngày \\
Slice & \texttt{WHERE} một chiều & năm 2017 \\
Dice & \texttt{WHERE} nhiều chiều & Sudeste $\times$ Home \& Furniture $\times$ Q4 \\
Pivot & \texttt{FILTER} & vùng $\times$ năm \\
\bottomrule\end{tabular}}\end{center}

\note{[S10] Mỗi thao tác là một câu SQL ngắn; phần 5 sẽ cho kết quả thật.}
\end{frame}

\begin{frame}[noprogressbar,noframenumbering]{MOLAP tính sẵn khối đa chiều, ROLAP tính lúc truy vấn, HOLAP kết hợp: nhóm chọn ROLAP}
\begin{itemize}
  \item Lý do: PostgreSQL không có khối đa chiều, dữ liệu 112.650 dòng hàng trả lời dưới 1 giây, kho luôn cập nhật sau mỗi lần nạp
\end{itemize}
\begin{center}\adjustbox{max width=\linewidth,max height=0.5\textheight}{\scriptsize\renewcommand{\arraystretch}{1.15}\begin{tabular}{llll}\toprule
\textbf{Tiêu chí} & \textbf{MOLAP} & \textbf{ROLAP} & \textbf{HOLAP} \\\midrule
Lưu trữ & khối đa chiều riêng & bảng quan hệ (lược đồ sao) & kết hợp hai cách \\
Tốc độ truy vấn & rất nhanh (tính sẵn) & phụ thuộc chỉ mục, materialized view & cân bằng \\
Độ tươi dữ liệu & phải xử lý lại khối & luôn tươi sau mỗi lần nạp & tùy phần \\
Công cụ & ví dụ SSAS & PostgreSQL + Power BI & -- \\
Nhóm chọn & -- & \ding{51} ROLAP & -- \\
\bottomrule\end{tabular}}\end{center}

\note{[S11] Đổi lại, truy vấn nặng được tăng tốc bằng materialized view; phần hiệu năng sẽ cho thấy lợi ích.}
\end{frame}

\begin{frame}[noprogressbar,noframenumbering]{BI đưa kho dữ liệu đến người ra quyết định qua mô hình ngữ nghĩa, KPI và dashboard}
\begin{itemize}
  \item Power BI: nạp CSV, quan hệ, measure DAX, dashboard
\end{itemize}
\khungchen[0.3\textheight]{CHÈN HÌNH / BẢNG: sơ đồ chuỗi dữ liệu $\rightarrow$ kho $\rightarrow$ mô hình $\rightarrow$ dashboard (vẽ)}

\note{[S12] Mọi KPI được định nghĩa một lần bằng DAX và đối chiếu với giá trị tính bằng SQL.}
\end{frame}
